\chapter{电磁超声换能与接触检测理论基础}

\section{电磁超声换能机理与接触边界建模}

电磁超声检测的独特价值在于，它将探头与结构之间的接触问题从传统压电超声所需的”声学耦合剂”转化为”电磁场、材料边界与机械约束是否共同产生可解释的回波”这一物理可观测性问题。在无人机平台上，这一点尤为关键：机体姿态、末端压力和探头升距会同时改变回波幅值与相位，使单一阈值不再可靠。本节将EMAT信号视为由激励、传播、反射和接收四个物理环节串联形成的完整可观测动态过程，从电磁-弹性耦合的第一性原理出发，依次推导洛伦兹力激励机制、弹性波传播规律，进而建立升距与回波质量之间的定量联系，最终将接触状态建模为多模态条件下的概率估计问题。

\subsection{涡流激励与洛伦兹体力}

在导电材料表面，EMAT探头中的交变线圈通入高频激励电流$I(t)$，在材料集肤深度内诱导出交变磁场，进而感生出涡电流。涡流与永磁体提供的静磁场$\mathbf{B}$相互作用，形成洛伦兹体积力$\mathbf{L}$：

\begin{equation}
\mathbf{L} = \mathbf{J} \times \mathbf{B}
\label{eq:lorentz_force}
\end{equation}

其中电流密度$\mathbf{J}$由材料电导率$\sigma$、局部电场$\mathbf{E}$和运动诱导项共同决定：

\begin{equation}
\mathbf{J} = \sigma\left(\mathbf{E} + \mathbf{v} \times \mathbf{B}\right)
\label{eq:current_density}
\end{equation}

式(\ref{eq:current_density})中的速度项$\mathbf{v} \times \mathbf{B}$直接体现了运动导体在静磁场中的感应效应。对于无人机检测而言，探头微小的法向摆动会改变升距与涡流分布，切向速度会改变等效运动项，因此回波变化不能简单归因于缺陷或厚度变化，而应在位姿与接触状态约束下联合解释。

式(\ref{eq:lorentz_force})所示的洛伦兹体力并非抽象的信号源，而是直接进入材料弹性动力学方程的体载荷项，决定超声波的初始能量、极化方向与空间分布。其强度正比于静磁场强度和涡流密度，而涡流密度又受升距（影响磁耦合效率）和激励频率（影响集肤深度）的共同调制。

\subsection{弹性动力学方程与波速}

超声波在材料内部的传播由弹性动力学方程描述。在各向同性线弹性近似下，以位移场$\mathbf{u}$为基本未知量，考虑洛伦兹体力$\mathbf{L}$作为外部激励源，运动方程可写为：

\begin{equation}
\rho \frac{\partial^2 \mathbf{u}}{\partial t^2} = (\lambda + \mu) \nabla(\nabla \cdot \mathbf{u}) + \mu \nabla^2 \mathbf{u} + \mathbf{L}
\label{eq:elastodynamic}
\end{equation}

其中$\rho$为材料密度，$\lambda$和$\mu$为Lam\'{e}常数。应力与应变满足线弹性本构关系$\bm{\sigma} = \lambda(\nabla\cdot\mathbf{u})\mathbf{I} + 2\mu\bm{\varepsilon}$。

对式(\ref{eq:elastodynamic})进行Helmholtz分解，可得纵波与横波的传播速度：

\begin{equation}
c_L = \sqrt{\frac{\lambda + 2\mu}{\rho}}, \quad c_T = \sqrt{\frac{\mu}{\rho}}
\label{eq:wave_speeds}
\end{equation}

式(\ref{eq:wave_speeds})为后续模型提供了两个重要先验约束：一是有效回波的到达时间应与材料厚度$d$处于同一物理量级，即$t_{\text{arrival}} \approx 2d / c_{L,T}$；二是接触状态变化引起的信号扰动应表现为连续的幅相变化，不应在相邻时间帧之间出现无法由无人机刚体动力学解释的突变。

\subsection{升距与耦合效率}

EMAT常被称为非接触式超声技术，但在无人机自主检测任务中，探头与被测表面的相对位置仍然是决定信号质量的核心变量。升距$h$为探头底部到被测表面的法向距离。升距过大时，涡流密度因磁耦合效率指数衰减而迅速降低；接触压力过大时，机体姿态扰动和表面摩擦又会破坏稳定扫查。

因此，本课题所说的接触状态并非传统压电探头意义上的“耦合剂接触”，而是一个由升距$h$、法向约束$F_n$、末端力学稳定性和回波可观测性共同定义的多因素工作状态。接触的成立需同时满足三个条件：（1）视觉上目标处于可接近区域，（2）位姿上法向误差和速度误差收敛至可接受区间，（3）超声上回波统计特征进入稳定区间。

\subsection{回波统计表征}

在时域上，接收回波可以近似写成带衰减包络和噪声项的调制信号：

\begin{equation}
s(t) = A(t) \cos(2\pi f_c t + \phi(t)) + n(t)
\label{eq:echo_model}
\end{equation}

其中包络幅值$A(t)$、中心频率附近的能量分布、相位延迟$\phi(t)$和噪声底$n(t)$共同构成接触状态的证据。接触边界改变时，等效声阻抗发生变化，反射系数$R$随之改变：

\begin{equation}
R = \frac{Z_2 - Z_1}{Z_2 + Z_1}
\label{eq:reflection}
\end{equation}

其中$Z_i = \rho_i c_i$为材料的声阻抗。当探头远离表面时，电磁能量向材料内部耦合不足，$A(t)$降低且波形不稳定；当探头接近有效工作区间时，回波能量增强，主要峰值的到达时间更稳定，短时统计特征呈现更好的重复性。

上述机理意味着接触判别不能仅依赖单帧幅值，而应综合观察时间窗口$[t-\tau, t]$内的多维度特征向量：

\begin{equation}
\Phi(t) = \left[E(t),\, A_{\max}(t),\, t_{\text{arrival}}(t),\, f_{\text{centroid}}(t),\, \kappa(t),\, \langle s_i, s_j \rangle_{\tau}\right]
\label{eq:feature_vector}
\end{equation}

其中$E(t)$为短时能量，$A_{\max}(t)$为峰值幅值，$t_{\text{arrival}}(t)$为到达时间，$f_{\text{centroid}}(t)$为频谱重心，$\kappa(t)$为峰度，$\langle s_i, s_j \rangle_{\tau}$为窗口内相邻帧的互相关。

\subsection{接触状态的概率建模}

本文将接触状态$c_t \in \{0, 1\}$建模为隐变量，并以回波、位姿和视觉特征为观测量。接触概率的估计可表述为条件概率形式：

\begin{equation}
P(c_t = 1 \mid \mathbf{Z}_t, \mathbf{P}_t, \mathbf{V}_t)
\label{eq:contact_prob}
\end{equation}

其中$\mathbf{Z}_t$为EMAT回波观测向量，$\mathbf{P}_t$为位姿误差观测向量，$\mathbf{V}_t$为视觉ROI特征向量。式(\ref{eq:contact_prob})的形式旨在明确“接触”应被看作概率估计问题而非确定性分类问题。概率输出能够表达临界接触、轻微脱离和稳定接触之间的连续过渡，并可直接接入飞行控制状态机，用于从目标接近、接触确认到保持扫查的平滑切换。

更进一步，接触边界的判别应具有滞回意识。无人机从远场接近表面时，回波增强通常伴随视觉深度减小和法向速度收敛；而从稳定接触转向脱离时，回波衰减往往先于视觉目标丢失出现。为此，需在时间窗口内同时建模“进入接触区”和“离开接触区”两类方向性证据，使判别结果不仅准确，而且符合运动过程的先后顺序。

\section{本章小结}

本章从电磁-弹性耦合第一性原理出发，系统地阐述了EMAT换能的物理机制与接触边界建模方法。推导了洛伦兹体力作为弹性动力学体载荷项的基本表达式，揭示了电流密度中速度诱导项与无人机姿态运动之间的物理联系；阐明了弹性波传播速度与材料参数的定量关系，为回波到达时间的先验约束提供了理论依据。在此基础上，建立了升距与耦合效率之间的联系，对回波信号进行了多维度统计表征，并将接触状态从确定性二值标签推广为概率化的多模态条件估计问题，指出了滞回判别在控制闭环中的必要性。本章的理论框架为后续章节的多模态融合算法设计和实验验证奠定了物理基础。

% ══════════════════════════════════════════════════════════════════
% 第三章 实验平台
% ══════════════════════════════════════════════════════════════════